% !TEX program = xelatex
% !TeX TS-program = xelatex
% Minimalist UNNC Beamer — compile with XeLaTeX (latexmk -xelatex main.tex)
% 画幅：aspectratio=169（16:9）
\documentclass[10pt,aspectratio=169,hyperref={colorlinks,citecolor=unncmid,urlcolor=unncmid,linkcolor=}]{beamer}
\usetheme{UNNC}
\usefonttheme{serif}
\usepackage{lipsum}
% 中文支持：需用 XeLaTeX 编译（latexmk -xelatex main.tex）
% 中文衬线宋体：macOS 用 Songti SC；其他环境（Linux/Overleaf 等）自动回退到
% TeX Live 自带的 FandolSong（按文件名经 kpathsea 加载，不依赖系统字体）
\usepackage[fontset=none, scheme=plain]{ctex}
\IfFontExistsTF{Songti SC}{%
  \setCJKmainfont{Songti SC}[AutoFakeSlant=0.2]%
  \setCJKsansfont{Songti SC}[AutoFakeSlant=0.2]%
  \setCJKmonofont{Songti SC}%
}{%
  \IfFontExistsTF{Noto Serif CJK SC}{%
    \setCJKmainfont{Noto Serif CJK SC}[AutoFakeSlant=0.2]%
    \setCJKsansfont{Noto Serif CJK SC}[AutoFakeSlant=0.2]%
    \setCJKmonofont{Noto Serif CJK SC}%
  }{%
    \setCJKmainfont{FandolSong-Regular}[Extension=.otf,BoldFont=FandolSong-Bold,AutoFakeSlant=0.2]%
    \setCJKsansfont{FandolSong-Regular}[Extension=.otf,BoldFont=FandolSong-Bold,AutoFakeSlant=0.2]%
    \setCJKmonofont{FandolFang-Regular}[Extension=.otf,BoldFont=FandolFang-Regular]%
  }%
}
\usepackage{hyperref}
% 英文衬线：XCharter（fontspec + Extension=.otf，Overleaf/TeX Live 均可用）
\setmainfont{XCharter}[
  Extension=.otf,
  UprightFont=*-Roman,
  BoldFont=*-Bold,
  ItalicFont=*-Italic,
  BoldItalicFont=*-BoldItalic,
]
% 等宽字体：优先 Maple Mono NF，其次 Maple Mono，未安装则回退默认
\IfFontExistsTF{Maple Mono NF}{%
  \setmonofont{Maple Mono NF}[Scale=MatchLowercase,Ligatures=TeX]%
}{%
  \IfFontExistsTF{Maple Mono}{%
    \setmonofont{Maple Mono}[Scale=MatchLowercase,Ligatures=TeX]%
  }{}%
}
% other packages
\usepackage{latexsym,amsmath,xcolor,multicol,booktabs}
\usepackage{amssymb}
\usepackage{graphicx}
\usepackage{tikz}
\usepackage{bm}
\usepackage{natbib}
\usepackage{wrapfig}
\usepackage{amsfonts}
\usepackage{ragged2e}
\usepackage{parskip}

\apptocmd{\frame}{}{\justifying}{} % Allow optional arguments after frame.

% 引用整体着色为 unncmid（与 definition / link 色一致）
\makeatletter
\NewCommandCopy{\unnc@citep@orig}{\citep}
\NewCommandCopy{\unnc@citet@orig}{\citet}
\RenewDocumentCommand{\citep}{o o m}{%
  \begingroup\color{unncmid}%
  \IfNoValueTF{#1}{\unnc@citep@orig{#3}}%
  {\IfNoValueTF{#2}{\unnc@citep@orig[][#1]{#3}}%
  {\unnc@citep@orig[#1][#2]{#3}}}%
  \endgroup}
\RenewDocumentCommand{\citet}{o o m}{%
  \begingroup\color{unncmid}%
  \IfNoValueTF{#1}{\unnc@citet@orig{#3}}%
  {\IfNoValueTF{#2}{\unnc@citet@orig[][#1]{#3}}%
  {\unnc@citet@orig[#1][#2]{#3}}}%
  \endgroup}
\makeatother

\newcommand{\theHalgorithm}{\arabic{algorithm}}
\theoremstyle{plain}
\newtheorem{axiom}{Axiom}
\newtheorem{claim}[axiom]{Claim}
\newtheorem{assumption}{Assumption}
\newtheorem{remark}{Remark}
\newtheorem{proposition}{Proposition}
\setbeamertemplate{theorems}[numbered]

% change for your title page information
\author[Author]{Your Name \; 作者姓名\\{\small\href{mailto:your.email@example.com}{\texttt{your.email@example.com}}}}
\title[Minimalist UNNC Beamer]{Minimalist UNNC Beamer}
\subtitle{A lightweight UNNC-style Beamer template with Chinese support}
\institute{University of Nottingham Ningbo China \\ 宁波诺丁汉大学}
\date{2026}

\def\cmd#1{\texttt{\color{unncmid}\footnotesize $\backslash$#1}}
\def\env#1{\texttt{\color{unncblue}\footnotesize #1}}
\definecolor{deepblue}{rgb}{0,0,0.5}
\definecolor{deepred}{rgb}{0.6,0,0}
\definecolor{deepgreen}{rgb}{0,0.5,0}
\definecolor{halfgray}{gray}{0.55}

\begin{document}

{
\setbeamertemplate{logo}{}
\begin{frame}
  \titlepage
  \begin{center}
    % Vector lockup (official UoN SVG → PDF). PNG kept in figures/ for raster use.
    % Slightly smaller, with space underneath, so the mark clears the footline.
    \includegraphics[width=0.31\linewidth]{figures/unnc_logo_full.pdf}\par
    \vspace{0.42cm}
  \end{center}
\end{frame}
}

\section{Introduction}

% 中文排版示例（可删除）：中英文混排、粗体、斜体
\begin{frame}{中文支持示例}
宁波诺丁汉大学（University of Nottingham Ningbo China, UNNC）成立于 \textbf{2004} 年，是中国第一所中外合作大学。

\begin{itemize}
  \item 中文正文使用衬线\textbf{宋体}（Songti SC / FandolSong），与英文 XCharter 衬线字体搭配
  \item 等宽字体优先 \texttt{Maple Mono}：\texttt{0O Il 1l \{[(=+])\}}（未安装则回退默认）
  \item {\itshape 斜体中文通过伪斜体自动生成，}数学公式不受影响：$\int_0^\infty e^{-x^2}\,\mathrm{d}x = \frac{\sqrt{\pi}}{2}$
  \item 标点符号使用全角样式：“引号”、‘单引号’、破折号——省略号……
\end{itemize}

\begin{theorem}[中文定理示例]
  若函数 $f(x)$ 在闭区间 $[a,b]$ 上连续，则 $f(x)$ 在 $[a,b]$ 上必有最大值与最小值。
\end{theorem}
\end{frame}

\begin{frame}{How to cite in your introduction?}
\lipsum[1][1-3]

\begin{enumerate}
  \item \lipsum[1][1-2] \citep{langley00}
  \item \citet{mitchell80} \lipsum[1][3-4]
  \item \lipsum[1][5-6] \citep{kearns89}
\end{enumerate}

Important contents are \textcolor{unncblue}{\bf bold Nottingham blue} \citep{MachineLearningI}.
\lipsum[1][1-4]
\end{frame}

\begin{frame}{Wrapping text around figures}
  \begin{wrapfigure}{r}{0.32\textwidth}
    \begin{center}
      \includegraphics[width=0.28\textwidth]{figures/unnc_crest.png}
    \end{center}
    \caption{UNNC crest (demo placeholder)}
  \end{wrapfigure}

  \lipsum[1][1-4]

  \lipsum[1][5-10]

  \lipsum[2][1-3]
\end{frame}

\begin{frame}{Mathematical expressions}
  \lipsum[3][1-5]

  \lipsum[4][1-3]
  \begin{equation}\label{eq:1}
  \bm Y = \bm G \bm X + \bm E,
  \end{equation}
  where
  \begin{itemize}
    \item $\bm G\in \mathbb{R}^{N\times SO}$ is a known gain matrix,
    \item $\bm E$ is the IID Gaussian error with $e_{ij}\sim\mathcal{N}(0,\sigma^2)$,
    \item \lipsum[2][1]
  \end{itemize}
\end{frame}

\section{Motivations}

\begin{frame}{Add block in text}
\lipsum[3][1-3]

\lipsum[3][1-2]

\lipsum[3][1-4]

\begin{block}{What is Few-Shot Learning?}
  \lipsum[3][1-4]
\end{block}
\end{frame}

\section{Methods}

\begin{frame}{Definition, Proposition and Theorem}

\begin{definition}[some explanations]
\lipsum[3][1-4]
\end{definition}

\begin{proposition}\label{prop:1}
\lipsum[3][5-8]
\end{proposition}

\begin{theorem}\label{theorem:generalization bound}
\lipsum[3][9-12]
\end{theorem}

\end{frame}

\section{Simulation}

\begin{frame}{Tables}
  \begin{table}
    \caption{Units for Magnetic Properties}
    \label{table}
    \setlength{\parskip}{0pt}%
    \footnotesize
    \setlength{\tabcolsep}{3pt}
    \begin{tabular}{p{40pt}p{100pt}p{140pt}}
      \toprule
      {\bfseries Symbol}&
      {\bfseries Quantity}&
      {\bfseries Conversion from Gaussian and \par CGS EMU to SI $^{\mathrm{a}}$} \\
      \midrule
      $\Phi $&
      magnetic flux&
      1 Mx $\to 10^{-8}$ Wb $= 10^{-8}$ V$\cdot $s \\
      $B$&
      magnetic flux density, \par magnetic induction&
      1 G $\to 10^{-4}$ T $= 10^{-4}$ Wb/m$^{2}$ \\
      $H$&
      magnetic field strength&
      1 Oe $\to 10^{3}/(4\pi )$ A/m \\
      $m$&
      magnetic moment&
      1 erg/G $=$ 1 emu \par $\to 10^{-3}$ A$\cdot $m$^{2} = 10^{-3}$ J/T \\
      $M$&
      magnetization&
      1 erg/(G$\cdot $cm$^{3}) =$ 1 emu/cm$^{3}$ \par $\to 10^{3}$ A/m \\
      4$\pi M$&
      magnetization&
      1 G $\to 10^{3}/(4\pi )$ A/m \\
      $\sigma $&
      specific magnetization&
      1 erg/(G$\cdot $g) $=$ 1 emu/g $\to $ 1 A$\cdot $m$^{2}$/kg \\
      \bottomrule
      \multicolumn{3}{p{280pt}}{Vertical lines are optional in tables. Statements that serve as captions for
        the entire table do not need footnote letters. }
    \end{tabular}
    \label{tab1}
  \end{table}
\end{frame}

\begin{frame}{References}%[allowframebreaks] if too many references
  \renewcommand*{\bibfont}{\footnotesize}
  \bibliography{ref}
  \bibliographystyle{apalike}
\end{frame}

\begin{frame}[plain]
  \begin{center}
    {\Huge\bfseries\color{unncblue} Thanks!}
    \vspace{1.2em}

    {\large\color{unncgrey} University of Nottingham Ningbo China}\\[0.4em]
    {\normalsize\color{unncgrey} 宁波诺丁汉大学}
  \end{center}
\end{frame}

\end{document}
